%%%%%%%%%%%%%%%%%%%%%%%%%%%%%%%%%%%%%%%%%%%%%%%%%%%%%%%%%%%%
\section{Literaturverzeichnisse}%
\label{sec:Literaturverzeichnisse}%
\index{Bibliografie}\index{Literaturverzeichnis}
%%%%%%%%%%%%%%%%%%%%%%%%%%%%%%%%%%%%%%%%%%%%%%%%%%%%%%%%%%%%
%
Defaultmäßig befinden sich alle Literaturangaben in der Datei
\printfilepathquoted{bib/Diss.bib}.
Falls ein davon abweichender Pfad oder zusätzliche Literaturdatenbanken
verwendet werden sollen, können entsprechende Pfade in der Datei
\printfilepathquoted{preambel/AllePfade.tex} mit
\lstinline|\addbibresource{...}|
festgelegt werden.
Bitte beachten, dass bei den Literaturdatenbanken kein \gls{gls:bibtex}-,
sondern \gls{gls:biblatex}-Format verwendet werden soll,
also Umlaute direkt eingegeben werden können.

Die Vorlage ermöglicht es, neben dem Haupt"=Literaturverzeichnis
jeweils separat durchnummerierte Listen der eigenen Publikationen und Patente
sowie der betreuten studentischen Arbeiten zu erstellen.
Standardmäßig werden alle Literaturangaben aus der Hauptliteraturdatenbank
\printfilepathquoted{bib/Diss.bib} extrahiert.
Durch Änderung der Definitionen der Pfade
\printkeyword{bibpathOwnPatents},
\printkeyword{bibpathOwnPublications} und
\printkeyword{bibpathStudentTheses}
in der Datei
\printfilepathquoted{preambel/AllePfade.tex}
können jedoch für die einzelnen Auflistungen 
jeweils eigene, speziell dafür angelegte Datenbanken verwendet werden.

Die Einbindung der Literaturverzeichnisse in das Manuskript erfolgt in der Datei
\printfilepathquoted{content/Inhalt-BackMatter.tex}.
Die Zusatz"=Literaturverzeichnisse werden separat erstellt und sind vom Haupt"=Literaturverzeichnis unabhängig.
Dazu werden in den entsprechenden Abschnitten der Datei,
welche jeweils mit Hilfe der \printkeyword{refsection}"=Umgebung abgegrenzt sind,
die kompletten Inhalte der jeweiligen Literaturdatenbank unsichtbar referenziert.
Dies geschieht mit dem \lstinline|\nocite{*}|"=Befehl.
Alternativ dazu können die einzelnen Werke mit
\lstinline|\nocite{Quelle1,Quelle2,...,QuelleN}|
explizit referenziert werden.

Sofern man eine einzige Datenbank für sämtliche Literaturverzeichnisse verwendet
und die unsichtbare Referenzierung jeweils mit dem \lstinline|\nocite{*}|"=Befehl erfolgt,
müssen die relevanten Referenzen für das jeweilige Zusatz"=Literaturverzeichnis
gefiltert werden.
Die Filterung der eigenen Publikationen, Patente und betreuten Arbeiten
für die einzelnen Verzeichnisse geschieht jeweils mit der Angabe
\lstinline|keyword=ownpubl|,
\lstinline|keyword=ownpatent|
oder
\lstinline|keyword=supervised|
in der Optionenliste des jeweiligen
\lstinline|printbibliography|"=Befehls.
Dazu müssen die Literatureinträge in der Literaturdatenbank freilich
entsprechend jeweils mit der Angabe
\lstinline|keywords = {ownpubl}|,
\lstinline|keywords = {ownpatent}|
oder
\lstinline|keywords = {supervised}|
versehen werden.

Bei Verwendung separater Literaturdatenbanken für die einzelnen Zusatzverzeichnisse
oder bei einer expliziten Aufzählung der einzelnen Veröffentlichungen mit dem
\lstinline|\nocite{Quelle1,Quelle2,...,QuelleN}|"=Befehl,
ist eine Filterung der Literaturdatenbank nach einem Keyword nicht notwendig.
Damit könnte die Notwendigkeit von Zusatzangaben in der Literaturdatenbank entfallen.
Allerdings muss dann auch die standardmäßig gesetzte Filter-Option in den jeweiligen
\lstinline|printbibliography|"=Befehlen
in der Datei  \printfilepathquoted{content/Inhalt-BackMatter.tex}
deaktiviert (sprich auskommentiert) werden.

Die Zusatzauflistungen sind von den Zitierungen im Text unabhängig.
Damit ergibt sich für eigene Veröffentlichungen,
die im Text explizit zitiert werden, eine doppelte Listung --
einmal in der Hauptliteraturliste (mit einer alphanumerischen Referenz)
und einmal in der durchnummerierten Liste eigener Publikationen.
Dies ist kein Bug, sondern ein Feature der Vorlage.

Bei der Auflistung eigener Veröffentlichungen und Patente ist 
in der Vorlage jeweils ein kleines Vorwort vorgesehen.
Darin können beispielsweise Zusatzhinweise zur Abgrenzung der Publikationen \bzw
der Patente, die während der Promotionszeit entstanden sind,
oder Hinweise zur Eingruppierung der Veröffentlichungen platziert werden.
Diese einleitenden Texte sind ebenfalls in der Datei
\printfilepathquoted{content/Inhalt-BackMatter.tex}
definiert und können dort editiert werden.

Bei der Erwähnung der einzelnen eigenen Veröffentlichungen und Patente in den Vorworten 
zu den jeweiligen Listen
muss die Einbindung mit Hilfe des Makros
\lstinline|\citefield{...}[labelnumberwidth]{labelnumber}|
statt der üblichen Einbindung mit dem Makro
\lstinline|\cite{...}|
erfolgen, damit die zu den Listen passende (nummerische) Formatierung verwendet wird.

Sind keine Vorworte vor den Listen eigener Veröffentlichungen und Patente erwünscht,
lassen sich diese durch die Auskommentierung der Optionen
\lstinline|prenote=OwnPublPrenote|
\bzw
\lstinline|prenote=OwnPatentPrenote|
in den jeweiligen 
\lstinline|printbibliography|"=Befehlen komplett ausblenden.